% =====================================================================
% main.tex -- DCC 2027 submission, dictforge
%
% RESTRUCTURED 2026-08-31 to match docs/DCC_STYLE.md (evidence-based
% format derived from 13 accepted DCC papers): 10 pages total including
% references (no appendix), ~4,500-4,800 words, section order
% Intro/Related/Background/Diagnosis/Method/Eval-Methodology/Results/
% Conclusion, contributions in prose, no Discussion/Limitations heading.
% Numbers affected by the in-flight refit-on-full fix (KNOWN_ISSUES.md
% #1) are wrapped in \pending{} -- see the restructuring report for the
% full list; do not resolve them here, they update when the fix lands.
%
% DCC "Information for Authors" formatting (re-verify against the actual
% DCC 2027 call before submission): regular papers 10 pages MAXIMUM
% including references and figures; single-column, 12pt; text area
% approximately 6in x 9in (docs/DCC_STYLE.md, verified against the
% 2025/2026 programs). No official .cls/.sty is on hand, so this uses
% `article` + `geometry` configured to match that text area.
% =====================================================================
\documentclass[12pt]{article}

\usepackage[letterpaper,left=0.85in,right=0.85in,top=0.85in,bottom=0.85in]{geometry}
\usepackage[utf8]{inputenc}
\usepackage[T1]{fontenc}
\usepackage{amsmath}
\usepackage{amssymb}
\usepackage{mathtools}
\usepackage{graphicx}
\usepackage{booktabs}
\usepackage{multirow}
\usepackage{xcolor}
\usepackage{url}
\usepackage{hyperref}
\hypersetup{colorlinks=true, linkcolor=black, citecolor=black, urlcolor=blue}
\usepackage{titlesec}
\usepackage{enumitem}
\usepackage{multicol}

% Tighten float/caption spacing -- page-budget measure, not a content
% change; DCC's own style packs floats more densely than article's
% defaults.
\setlength{\intextsep}{6pt plus 2pt minus 2pt}
\setlength{\textfloatsep}{8pt plus 2pt minus 2pt}
\setlength{\abovecaptionskip}{4pt}
\setlength{\belowcaptionskip}{2pt}
\emergencystretch=1.5em
\renewcommand{\arraystretch}{0.82}

% Squeeze inter-entry spacing in the bibliography -- page-budget measure
% only, does not change font size or entry content.
\let\oldthebibliography\thebibliography
\renewcommand{\thebibliography}[1]{%
  \oldthebibliography{#1}%
  \setlength{\itemsep}{0pt}%
  \setlength{\parskip}{0pt}%
}

% Pack floats more tightly against text rather than deferring them to
% isolated float pages -- a page-budget measure, not a content change.
\renewcommand{\topfraction}{0.9}
\renewcommand{\textfraction}{0.05}
\renewcommand{\floatpagefraction}{0.6}
\setcounter{topnumber}{4}
\setcounter{totalnumber}{6}

% Tighter section/list spacing -- DCC's own style is more compact than
% article's defaults; typographic only, no content change.
\titlespacing*{\section}{0pt}{1.4ex plus 1ex minus .2ex}{0.8ex plus .2ex}
\titlespacing*{\subsection}{0pt}{1.1ex plus 1ex minus .2ex}{0.6ex plus .2ex}
\setlist{itemsep=1pt, topsep=2pt, parsep=0pt}
\setlength{\abovedisplayskip}{4pt plus 2pt minus 2pt}
\setlength{\belowdisplayskip}{4pt plus 2pt minus 2pt}

% ---------------------------------------------------------------------
% \pending{...}: greppable marker for numbers known to be in flux ahead
% of the refit-on-full fix (KNOWN_ISSUES.md #1) and the corresponding
% equal-compute recomputation (results/corrections_v2.csv). Renders the
% wrapped text unchanged; sweep with `grep '\\pending{' main.tex` once
% corrected numbers land, then delete the macro.
% ---------------------------------------------------------------------
\newcommand{\pending}[1]{#1}

% Lightweight proposition environment for the trainer's two guarantees.
\newtheorem{proposition}{Proposition}

\title{Dictionary Training for Zstandard, Revisited:\\ Diagnosis and a Level-Aware Trainer}
% TODO: placeholder affiliation -- fill in real institution/affiliation before submission.
\author{
  Yugen \\
  \textit{TODO: affiliation} \\
  \texttt{19thkingisreal@gmail.com}
}
\date{}

\begin{document}
\maketitle

% =====================================================================
\begin{abstract}
Zstandard's dictionary trainer (COVER/fastCOVER)~\cite{liao2016cover} is
the de facto standard for small-record compression. However, it has
received no algorithmic attention since 2018: requesting a larger
dictionary can silently yield a smaller one, and construction ignores
size and level. Existing tools reuse the same heuristic and a fixed
objective that cannot tell an efficient dictionary from a starved one.
This paper presents a validation-driven, level-aware trainer that
searches size and trainer family under an anytime never-worse
guarantee, then refines content and seeds offsets at high levels on
five corpora. Against the deployed \texttt{zstd
--train} default, our trainer improves compression ratio by
\pending{3.1}--14.3\% at level 3 and 11.5--32.1\% at level 19.
Matched-size rebuilds show up to 6.9\% of the level-19 gain survives at
identical budgets, versus 0--11\% at level 3.
\end{abstract}

% =====================================================================
\section{Introduction}
\label{sec:introduction}

Dictionaries are how zstd compresses small data, and small data is
where modern systems live: key-value stores, logs, telemetry, and the
emerging shared-dictionary web transport~\cite{rfc9842}. One trainer
serves all of it: COVER and its fast approximation fastCOVER~\cite{liao2016cover},
a 2016 heuristic frozen since 2018, whose own maintainer has stated
that the ideal dictionary builder ``does not exist.''

Yet the trainer has documented pathological failures. A production
engineer reported that raising \texttt{--maxdict} from 512\,KB to
550\,KB collapsed compression ratio from $90\times$ to $14\times$ on
their data~\cite{zstdissue4127}; the thread was closed without an
explanation, and the trainer's behavior is unchanged in the current
release. We show this signature is reachable on zstd's own benchmark
corpus and trace it to a specific selection-objective defect, alongside
two further defects: the trainer never searches dictionary size, and it
optimizes construction decisions for a single compression level
regardless of which level the caller actually uses.

Guided by this diagnosis, we build a trainer that treats dictionary
construction as a search adjudicated by measurement rather than by
proxy, with every stage gated on held-out validation rather than on the
training data it searched over.

Our contributions are as follows. We give the first mechanistic,
quantitative diagnosis of the incumbent trainer, tracing budget
degeneracy to its selection objective, quantifying size blindness, and
confirming level blindness three independent ways; the budget-degeneracy
finding is consistent with a sensitivity collapse reported upstream but
never explained~\cite{zstdissue4127}. We build a level-aware,
validation-driven trainer with an anytime-safety guarantee
(Propositions~\ref{prop:budget-monotonicity} and~\ref{prop:anytime-floor}),
each stage individually gated and validation-checked. We evaluate it on
a 5-corpus, 2-level benchmark against the deployed default, a klauspost
\texttt{builddict}~\cite{klauspost} baseline, an equal-wall-clock stock
sweep, a no-training control, and matched-size rebuilds at Cassandra's,
RocksDB's, and (a setting motivated by) ScyllaDB's production budgets,
isolating how much of the headline gain is size versus method. The
resulting tool emits 100\%-standard-format dictionaries, requiring no
change to any zstd library or client, and is open source together with
all experiments.

% =====================================================================
\section{Related Work}
\label{sec:related-work}

\textbf{Dictionary construction.} The trainer descends from Liao,
Petri, Moffat, and Wirth's relative Lempel-Ziv dictionary
construction~\cite{liao2016cover}, a covering problem over sampled
substrings; zstd adopted it as \texttt{--train-cover} and later added
the fastCOVER approximation that is today's CLI default, unchanged
since 2018 and never described in a paper of its own. We reuse the
covering heuristic unmodified; our contribution is the selection layer
around it. Level-dependence is not a new observation -- a zstd
maintainer noted in 2019 that training at the level you actually
compress at ``will tune the dictionary better''~\cite{zstdissue1572},
and a 2022 request to expose a level parameter was declined for ABI
stability~\cite{zstdissue3213} -- but the \emph{mechanism} (three
independent, opposite-sign inversions traced to hash-table overwrite
vs.\ full indexing, \S\ref{sec:diagnosis-level-blindness}) and
\emph{magnitude}, including the first measurement of the disabled
repcode-seeding path, are new. \texttt{klauspost/compress}~\cite{klauspost}
is the only other tool we know of that emits standard-format zstd
dictionaries and exposes a level-targeting flag; we benchmark it
directly and it loses to the stock default on every cell we ran.

\textbf{RLZ, methodology, and theory.} Our refinement stage
(\S\ref{sec:design-stage2}) adapts RLZ reference pruning and
selection~\cite{hoobin2011rlz,puglisi2023rlz}: the novelty is doing it
against the production encoder's own sequence stream, at match
granularity, gated by measured compressed size. Optimal dictionary
choice is intractable in general -- NP-complete~\cite{storer1982textual},
APX-hard as smallest-grammar~\cite{charikar2005smallestgrammar,casel2019grammar}
-- so our propositions (\S\ref{sec:design-guarantees}) deliberately
constrain only the \emph{selection} rule, not the underlying trainer;
domain-specific compression~\cite{machete2024} and LZ77
acceleration~\cite{lzhv2025} are recent DCC neighbors in a different
corner of the space. Methodologically we follow
recent practical compression work that measures precisely why the
incumbent underperforms before designing against
it~\cite{boncz2020fsst,afroozeh2024alp}, reporting losses
(\S\ref{sec:results-parity}--\S\ref{sec:results-equalcompute}) as well
as wins.

\textbf{Deployment.} Meta's managed-compression practice, disclosed only
in a blog post~\cite{metaengineering2018zstd}, and the successor project
OpenZL~\cite{openzl}, treat per-workload specialization as known
industrial practice with little published methodology; we target the
specific, publicly reachable trainer instead. Deployment of trained
dictionaries varies sharply by system and path
(\S\ref{sec:background-deployment}). The web's shared-dictionary
transport~\cite{rfc9842} and Brotli's own dictionary-aware
format~\cite{alakuijala2019brotli} see modest uptake so far -- Chrome
usage sits at roughly 0.09\% of page
loads~\cite{chromestatususecounters}, and Cloudflare's implementation
remains a passthrough beta that itself generates no
dictionaries~\cite{cloudflaresharedddict} -- context for why
\S\ref{sec:results-parity} matches specific deployment configurations
rather than only our own chosen sizes.

% =====================================================================
\section{Background}
\label{sec:background}

\subsection{Dictionaries, indexing, and the trainer}
\label{sec:background-zstd-dicts}

A zstd dictionary is a self-describing binary artifact -- magic number,
ID, entropy tables (12 bytes of seeded recent-offset values), then
\texttt{Content}, a virtual ``past'' matches may reference by
offset~\cite{rfc8878}. Both loaders parse this independently and
symmetrically, so a dictionary from any compliant writer is a drop-in
replacement for \texttt{zstd --train}: \emph{fact 1}. \texttt{Content}
is not indexed uniformly, though: fast strategies hash it into one or
two tables where later positions can overwrite earlier ones, while
high-level binary-tree strategies retain every position -- how much of
``the dictionary'' is reachable is a function of level: \emph{fact 2},
the mechanism behind level blindness (\S\ref{sec:diagnosis-level-blindness}).

The trainer solves a covering problem: COVER~\cite{liao2016cover} builds
an exact suffix array over $d$-length substrings and scores each by
document frequency, deduplicated per sample (a dictionary only helps
the \emph{first} reference to a substring); fastCOVER trades this for a
lossy, streaming hash-table count. Both partition the corpus into
epochs, greedily pack the best-scoring $k$-byte window per epoch into
the output buffer back forward (strongest content nearest the data,
cheapest to reference), and grid-search $(d,k)$, keeping whichever
candidate achieves the smallest \emph{measured} check-set compressed
size -- selection is literal compression, not a proxy. \texttt{zstd
--train} with no flags, and \texttt{ZDICT\_trainFromBuffer} used by
every downstream integration we examine, both hard-code
fastCOVER-optimize at $d{=}8$, with a 4-step grid on $k$, a 0.75
train/check split, scored at level 3, \texttt{--maxdict} defaulting to
110\,KB; repeat-offset seeding is computed and discarded behind
\texttt{\#if 0}, and the size-response \texttt{shrinkDict} path is
hard-coded off. All claims here were checked at file:line precision
against \texttt{facebook/zstd main@82d322c} (v1.6.0).

\subsection{Deployment reality}
\label{sec:background-deployment}
\begin{sloppypar}
Production systems mostly call the same entry point, or,
with no trainer, \texttt{ZDICT\_finalizeDictionary}, which skips
content selection and attaches only the format header around
caller-supplied content. \textbf{RocksDB}'s dictionary path is opt-in and off by
default; once enabled, it uses our baseline trainer path by default, at
110\,KB, level 3. \textbf{ScyllaDB} trains RPC and SSTable dictionaries
through the same 110\,KiB call, but only one is deployed by default:
RPC training is off, so our level-1 cell is a setting \emph{motivated
by} that path, not a production claim; SSTable dictionary compression,
by contrast, has been that vendor's default compressor since release
2025.4, so the 110\,KiB/level-3 configuration we test as ``the deployed
default'' is production-representative there. \textbf{Cassandra}'s
CEP-54 is not merged or shipped -- its tracking epic is still in
progress -- so we keep its 64\,KiB code default as an evaluation
target, not a production claim. Each budget used in
\S\ref{sec:results-parity} is labeled by what it is: shipped default,
evaluation setting motivated by an off-by-default path, or unreleased
branch.
\end{sloppypar}

% =====================================================================
\section{Diagnosis (measurements on real corpora)}
\label{sec:diagnosis}

\subsection{Budget degeneracy (the ``cliff'')}
\label{sec:diagnosis-cliff}

Consider a user who has trained a dictionary with \texttt{zstd
--train-cover} and raises \texttt{--maxdict} for a better one. On
\texttt{github\_users} -- zstd's own published benchmark sample set --
this produces Figure~\ref{fig:f1-cliff}'s step function: at requests of
262--524\,KiB the trainer emits exactly what was asked and ratio holds
near 10; past 1\,MiB it emits only $\sim$415--425\,KiB regardless of how
much more is requested, and ratio falls to 7.49 (a 24\% loss) and then
6.95 (31\%) -- \emph{smaller output at a larger request}, with no
warning at any verbosity level. In August 2024 an engineer integrating
zstd dictionaries at Roblox reported the same signature upstream -- a
512\,KB budget yielding a $90\times$ ratio against 550\,KB yielding
$14\times$ on their data~\cite{zstdissue4127} -- in a thread closed
without the trainer's behavior changing. We do not reproduce that
magnitude, but the mechanism below accounts for the signature.

\textbf{Mechanism.} Three interacting behaviors, traced to source and
confirmed by instrumented rebuild against
\texttt{facebook/zstd main@82d322c}, produce the collapse.
\textbf{M1:} epoch count scales with the \emph{requested} budget, not
the corpus (\texttt{cover.c:734-749}) -- at $k{=}50$ and a 2\,MiB
request, a 36\,MB training set becomes $\sim$42{,}000 sub-kilobyte
epochs. \textbf{M2:} the build loop exits after a hard-coded run of dead
epochs (fastCOVER: 10, \texttt{fastcover.c:406}), trivially reached on
redundant corpora, leaving the output buffer substantially unfilled.
\textbf{M3 (root cause):} the selection objective scores each candidate
by its \emph{emitted} size plus check-set compressed size
(\texttt{cover.c:899,998}), unable to distinguish ``small because
efficient'' from ``small because starved''; past a threshold the
optimizer deterministically selects the most degenerate build --
confirmed on \texttt{gharchive} at 2\,MiB, where the winner had the
\emph{worst} compression term of all candidates and won purely on a
1.62\,MB size advantage from its own failed fill. Roughly two thirds of
the \texttt{github\_users} collapse traces to this selection error, one
third to genuine oversize dilution. A minor \textbf{M4} compounds it:
finalize shrink-to-fit keeps the \emph{first} bytes when trimming for
the entropy header (\texttt{zdict.c:902-935}), discarding the tail --
the highest-scoring content the builder placed nearest the data on
purpose. Nobody noticed because both trainers assign
\texttt{ctx->displayLevel} and then \texttt{memset} the context to zero
a few lines later, silently suppressing every build-loop diagnostic at
every \texttt{-v} level. M2, M3, and \texttt{displayLevel} are
implementation defects; content-exhaustion pinning is inherent;
\texttt{--maxdict} semantics are a documentation gap; we have drafted
four upstream fixes.

\begin{figure}[t]
  \centering
  \includegraphics[width=\linewidth,height=0.22\textheight,keepaspectratio]{../figures/f1_cliff.pdf}
  \caption{Requested dictionary size vs.\ emitted size and vs.\
    compression ratio, one panel per corpus, marking degenerate builds
    (F1, generated by \texttt{tools/make\_figures.py}).}
  \label{fig:f1-cliff}
\end{figure}

\subsection{Size blindness}
\label{sec:diagnosis-size-blindness}

Ratio-vs-budget curves from 2KB to 2MB (Figure~\ref{fig:f2-size-curves})
show peaks varying from 32KB to over 2MB depending on corpus and level;
the default 110KB budget loses up to 22.4\% (gharchive, level 19), and
the curves are locally non-monotonic. The trainer searches $(d,k)$ but
never size; the \texttt{shrinkDict} path that would let it respond to
size is dead code.

\begin{figure}[t]
  \centering
  \includegraphics[width=\linewidth,height=0.24\textheight,keepaspectratio]{../figures/f2_size_curves.pdf}
  \caption{Compression ratio vs.\ dictionary size, one panel per
    (corpus, level), with the 110KB default marked (F2, generated by
    \texttt{tools/make\_figures.py}).}
  \label{fig:f2-size-curves}
\end{figure}

\subsection{Level blindness (three independent confirmations)}
\label{sec:diagnosis-level-blindness}

The trainer optimizes for a single level, though construction decisions
invert across levels. \textbf{(a)} Size-response inverts by level:
levels 1--3 peak small then decline, while level 19's response rises
monotonically. \textbf{(b)} Repcode seeding: $+0.3$--$0.4\%$ at level
19, $-0.2$--$0.6\%$ at level 1 -- the first proper evaluation of the
\texttt{\#if 0} path. \textbf{(c)} Refinement: $+1.7$--$6.6\%$ at level
19, approximately 0 at level 3. Fast levels index only a suffix (hash
overwrite) with greedy parsers, while \texttt{btopt} indexes everything
and cost-models its choices; the same pattern should recur whenever the
consumer is a different codec entirely -- LZ4's match window, for
instance, is capped at 64\,KiB -- though we have not measured that
cross-codec case in the rebuilt environment.

% =====================================================================
\section{Method: a validation-driven, level-aware trainer}
\label{sec:design}

The diagnosis in \S\ref{sec:diagnosis} suggests a design rather than an
algorithm: the stock trainer's failures are failures of
\emph{selection}, not of its covering heuristic, which is sound and
which we retain unchanged. We treat dictionary construction as
blackbox optimization over a discrete candidate space -- minimizing
total compressed size of a corpus under a size-capped dictionary, an
NP-hardness heritage traceable to Storer--Szymanski
substitution~\cite{storer1982textual} and smallest-grammar
APX-hardness~\cite{charikar2005smallestgrammar,casel2019grammar} -- with
a $\sim$1-second exact oracle (a real compression run). Evaluation is
cheap, so the design question is how to spend an evaluation budget.

\textbf{Data discipline.} The caller supplies one training directory,
split 85/15 by seed into a \emph{fit} set (what any underlying trainer
sees) and a \emph{validation} set (what every accept/reject decision is
measured on); the user's held-out data is never opened. This bounds
overfitting in our greedy per-stage loops -- the largest fit-to-heldout
deviation observed is 0.14\%.

\subsection{Stage 1: size and family search}
\label{sec:design-stage1}

Stage 1 sweeps a geometric size ladder from 2\,KiB to the caller's cap
with both fastCOVER and COVER, keeping whichever candidate scores best
on validation at the target level. The \emph{floor candidate} --
fastCOVER at $\min(\text{cap}, 110\,\text{KiB})$, precisely what
\texttt{zstd --train} would produce -- runs first unconditionally, so
every later decision is a comparison against what the user already had;
this makes the tool never-worse by construction
(Proposition~\ref{prop:anytime-floor}). FastCOVER sizes run before
COVER sizes (\emph{cheap-first ordering}), so a truncated time budget
cuts the expensive tail, not size coverage -- a size-first ordering
instead let large corpora exhaust the budget on COVER at small sizes,
losing to the default by up to 19\%. A \emph{degeneracy guard} flags
emitted size against requested size, but since selection is by measured
compression, not requested size (\S\ref{sec:diagnosis-cliff}, M3), a
starved build cannot win by being small. A ladder winner is
\textbf{refit on the full training directory} and kept only if
validation does not regress -- without this the tool loses $\sim$0.1\%
to the default whenever no ladder entry beats the floor.

\subsection{Stages 2 \& 3 (target level $\ge 16$): refinement and repcode seeding}
\label{sec:design-stage2}

Stage 1 chooses \emph{which} dictionary; Stage 2 asks what the covering
heuristic never asks: when real compressions run against it, which
bytes do they reference? We extract the sequence stream from fit-set
compressions, map matches to the regions they consumed, evict regions
never touched or in the weakest referenced decile, and refill from the
worst-compressing fit files, compacting proven content toward the tail
(\S\ref{sec:diagnosis-cliff}, M4). Each round is kept only if it
improves fit-set compressed size, stopping at the first rejection, and
the offset-to-region mapping self-tests against a planted synthetic
region before the stage may run at all. Literally-dead content is rare
(0--4\%), so gains come from displacing \emph{weakly} referenced
content. Stage 3 seeds the three recent-offset values the encoder
starts each frame with, which the stock trainer computes and discards,
writing $\{1,4,8\}$: we measure the fit set's top three offsets and
patch them in, keeping the result only if validation improves. Both
stages gate on target level $\ge 16$ because fast levels index only
part of a large dictionary into small hash tables and parse greedily,
so curated content and seeded offsets cannot be exploited and can
actively mislead the parser, while level 19's optimal parser indexes
everything and cost-compares its options.

\subsection{Guarantees}
\label{sec:design-guarantees}

\begin{sloppypar}
Let $S$ be the size ladder, $B$ the caller's budget cap, and
$C(B) = \{(t,s) : t \in \{\text{fastCOVER}, \text{COVER}\},\ s \in S,\ s
\le B\} \cup \{\text{floor}(B)\}$ the Stage-1 candidate set, where
$\text{floor}(B) = (\text{fastCOVER}, \min(B, 110\,\text{KiB}))$. The
tool returns $\arg\max_{c \in C(B)} r_{\text{val}}(c)$, the candidate
with the best measured validation-split compression ratio.
\end{sloppypar}

\begin{proposition}[Budget monotonicity]
\label{prop:budget-monotonicity}
For $B_1 \le B_2$, $C(B_1) \subseteq C(B_2)$, so the selected validation
ratio is monotone non-decreasing in $B$.
\end{proposition}

\begin{proposition}[Anytime floor]
\label{prop:anytime-floor}
For any time budget sufficient to evaluate $\text{floor}(B)$ -- seconds,
and it is evaluated first -- the selected validation ratio is at least
that of the incumbent trainer at the same cap.
\end{proposition}

Both are immediate from the construction; their content is not
mathematical depth but the observation that the \emph{deployed} trainer
satisfies neither -- Proposition~\ref{prop:budget-monotonicity} fails
empirically by 24--31\% (\S\ref{sec:diagnosis-cliff}), and
Proposition~\ref{prop:anytime-floor} fails because there is no floor at
all, DCC-normal scope for a guarantee on a systems contribution. The
guarantees hold on the validation measure, not held-out data
(\S\ref{sec:results} quantifies the gap), and
Proposition~\ref{prop:budget-monotonicity} constrains the
\emph{selected} candidate, not the underlying trainers -- a starved
build is still a starved build, we merely decline to select it.

% =====================================================================
\section{Evaluation Methodology}
\label{sec:evaluation-methodology}

Five small-record corpora spanning deployment shapes, including zstd's
own benchmark set: \texttt{github\_users} (9{,}114 JSON user records),
\texttt{gharchive} (20{,}000 GitHub events), \texttt{weblogs} (15{,}000
batches of 20 NASA-HTTP log lines), \texttt{apijson} (10{,}000 Crossref
records), \texttt{csvrows} (10{,}000 50-row NYC 311 CSV fragments), each
split 80/20 by seeded shuffle, with heldout data compressed exactly
once per configuration. We evaluate at level 3 (RocksDB/ScyllaDB
SSTable/Cassandra default) and 19 (archival), plus level 1 for a cell
motivated by ScyllaDB's off-by-default RPC path
(\S\ref{sec:background-deployment}), on an Apple M4 against zstd
\texttt{82d322c} (v1.6.0), every dictionary round-trip-verified against
an independent stock binary; x86 replication of throughput is planned
but not run.

We compare against five references: the deployed default (\texttt{zstd
--train} at 110\,KiB), a no-training control (random samples
concatenated to the same budget), \texttt{klauspost/compress}'s
\texttt{builddict}~\cite{klauspost}, an equal-wall-clock sweep of the
stock optimizer, and matched-size/production rebuilds at Cassandra's
64\,KiB budget, RocksDB's opt-in no-trainer fallback, and the
ScyllaDB-motivated level-1 setting (\S\ref{sec:background-deployment}).
Table~\ref{tab:t1-main} reports the \emph{tuned} (Stage 1 only) and
\emph{full} (all stages) variants of our trainer against these
baselines.

Midway through this work the machine's temporary storage was cleared,
destroying the experimental tree; we rebuilt it from recorded
specifications -- re-fetching every corpus, recompiling, re-running
every experiment -- and report only rebuilt-environment numbers.
Headline behaviors, including the budget-degeneracy collapse,
reproduced on freshly downloaded data with a fresh split, an unusually
strong reproducibility check.

Full per-corpus tables for the equal-compute sweep and the
\texttt{klauspost} baseline, the benchmark corpora, and the trainer
itself are available at \texttt{<REPO-URL-TODO>} (commit
\texttt{<HASH-TODO>}).

% =====================================================================
\section{Results}
\label{sec:results}

Table~\ref{tab:t1-main} and Figure~\ref{fig:f4-gains} give the headline
comparison against the deployed default; \S\ref{sec:results-main}
states it plainly, \S\ref{sec:results-parity} isolates how much of it
is dictionary size versus method by rebuilding baselines at matched
sizes and at production-motivated budgets, and
\S\ref{sec:results-equalcompute} checks whether the gain survives when
the incumbent is given equal wall-clock time and compares against the
only other public dictionary trainer we found.

\subsection{Main results}
\label{sec:results-main}

\input{tables/t1_main.tex}

\begin{figure}[t]
  \centering
  \includegraphics[width=\linewidth,height=0.17\textheight,keepaspectratio]{../figures/f4_gains.png}
  \caption{Percent compression-ratio improvement of \emph{tuned} and
    \emph{full} over the deployed default, by corpus and level (F4,
    generated by \texttt{tools/make\_figures.py}).}
  \label{fig:f4-gains}
\end{figure}

Against the deployed default, our trainer improves compression ratio on
every corpus at both levels (Table~\ref{tab:t1-main}): between
\pending{+3.1\%} and +14.3\% at level 3 and between +11.5\% and +32.1\%
at level 19, averaging roughly \pending{+8\%} and +23\% respectively
(Figure~\ref{fig:f4-gains}). The no-training control -- random training
samples concatenated to the same 110\,KiB -- trails the stock trainer
everywhere, confirming the covering heuristic earns its keep. The
\emph{tuned} column is Stage 1 alone; \emph{full} adds Stages 2--3,
contributing +1.0 to +9.0 percentage points at level 19, positive on
all five corpora. At level 3, where Stages 2--3 are gated off,
\emph{tuned} and \emph{full} should coincide exactly; where they
instead differ, by \pending{0\% to -2.7\%} and most visibly on
\texttt{github\_users} (\pending{-2.53\%}), the cause is a validation-gate
contamination in the full-refit step (\S\ref{sec:design-stage1}) rather
than genuine stage-2/3 contribution -- a fix is in progress and every
\pending{} figure in this section will change once it lands.

\subsection{Isolating size from method}
\label{sec:results-parity}

\begin{figure}[t]
  \centering
  \includegraphics[width=\linewidth,height=0.24\textheight,keepaspectratio]{../figures/f3_pareto.png}
  \caption{Compression ratio vs.\ dictionary size across every
    trainer/size variant we ran, one panel per (corpus, level); our
    \emph{tuned} and \emph{full} variants sit on or near the Pareto
    frontier in every panel (F3, generated by
    \texttt{tools/make\_figures.py}).}
  \label{fig:f3-pareto}
\end{figure}

Our trainer chooses its own size, usually larger than the 110\,KiB
default, so we rebuilt the baseline at \emph{our} winning size for
every corpus and level (Table~\ref{tab:t4-parity},
Figure~\ref{fig:f3-pareto}). At level 3, almost none of the headline
advantage survives -- two corpora match exactly, two retain
+0.2\%/+1.5\%, so 0--11\% of the gain is method rather than size. At
level 19 a real margin remains: +6.2\%/+6.9\%/+2.8\%
(\texttt{github\_users}/\texttt{weblogs}/\texttt{csvrows}) at identical
budgets. This defines the contribution: at fast levels the trainer
\emph{finds} the right size safely and automatically, a size the stock
trainer could in principle also have produced; at high levels
refinement and seeding add compression no choice of budget recovers.
Three of ten matched-size cells could not be run as intended -- the
stock trainer, asked for our winning budget, itself emitted
substantially less -- so the size confound is not fully separable from
\S\ref{sec:diagnosis-cliff}.

\input{tables/t4_parity.tex}

Reproducing Cassandra's shipped default and two further
deployment-motivated settings (\S\ref{sec:background-deployment}): at
Cassandra's 64\,KiB level-3 budget we beat the stock trainer on three
corpora (+0.4\% to +1.7\%) and tie exactly on two, no losses -- ties
occur where no ladder candidate beats the incumbent-equivalent floor,
the anytime guarantee (\S\ref{sec:design-guarantees}) operating as
designed. At a level-1, 110\,KiB setting motivated by ScyllaDB's
(off-by-default) RPC path, we win on all five corpora, +1.6\% to
+8.1\%. RocksDB's opt-in, non-default no-trainer fallback ranks between
no dictionary and the stock trainer in nine of ten cells and trails our
pipeline by a median of 18\%.

\subsection{Equal compute and alternative trainers}
\label{sec:results-equalcompute}

Our trainer spends minutes to tens of minutes (41\,s--59\,min across
corpora and levels) where \texttt{zstd --train} spends 0.2--3.2\,s. The
corresponding manual control gives the
incumbent our wall-clock budget to sweep sizes, selecting by measured
validation compression: this closes the gap at level 3
(\pending{1.5\%} above to \pending{9.5\%} below our result, a wash)
while at level 19 we still win on \pending{four of five} corpora,
\pending{+1.6\%} to \pending{+8.5\%}, losing narrowly on
\texttt{apijson} (\pending{-1.1\%}) -- these equal-compute figures
predate the refit-on-full fix above and will be replaced once corrected
values are computed. The sweep also reproduced the budget-degeneracy
pathology -- several large-budget candidates emitted dictionaries
25--30\% worse than at half the budget, surviving only because our
harness selects by measured compression, not requested size
(\S\ref{sec:diagnosis-cliff}, M3) -- the clearest demonstration that the
safety property in \S\ref{sec:design-guarantees} matters.

\texttt{klauspost/compress}'s \texttt{builddict}~\cite{klauspost}, the
closest prior art to our level-awareness claim, loses to the stock zstd
default on all twenty cells we ran, emitting $\sim$17\,KiB on
\texttt{github\_users} regardless of a 1\,MiB request; we report this
without enthusiasm -- it should have been a harder baseline to beat.

All results above are single-machine (Apple M4) measurements from the
rebuilt post-wipe environment; throughput, bootstrap confidence
intervals, and cross-codec transfer to Brotli and LZ4 were designed but
not re-executed here, so we report none rather than reuse figures we
could not recompute.

% =====================================================================
\section{Conclusion}
\label{sec:conclusion}

Diagnosis-first beats algorithm-first: a $\sim$1-second exact oracle
changes the design space. A 2016 covering heuristic, unmodified, plus a
validation-driven selection layer around it yields double-digit
archival gains, gains that survive matched-size and equal-compute
scrutiny, and guaranteed-safe defaults at every budget -- while the
safe-automation story at fast levels is, on its own honest terms, still
worth shipping. Two things remain open: byte-identical reproducibility
across reruns, an assumption the safety guarantees lean on, was not
independently re-checked after the post-wipe rebuild, and x86 throughput
replication is still pending. We have drafted, but not yet upstreamed,
four fixes to the incumbent trainer itself -- charging its objective by
requested rather than emitted capacity, scaling its dead-epoch
tolerance, fixing the suppressed diagnostics, and shrinking dictionaries
from the front instead of the back. Sharing Stage 1's search across
levels, and a speed-constrained mode that trades ratio for a bounded
training budget, are the natural next steps given the training-time
cost measured in \S\ref{sec:results-equalcompute}.

% =====================================================================
% Two-column, footnote-size reference list -- a page-budget technique,
% not a content change; every entry is byte-identical to references.bib.
{\footnotesize
\begin{multicols}{2}
\bibliographystyle{abbrv}
\bibliography{references}
\end{multicols}
}

\end{document}

% =====================================================================
% TODO tracker (pre-submission), carried over verbatim -- not part of
% the rendered paper.
% =====================================================================
% - [ ] Sweep every \pending{...} once the refit-on-full fix
%   (KNOWN_ISSUES.md #1) and the corrected equal-compute sweep
%   (results/corrections_v2.csv, variant=default_sweep_winner) land; see
%   the restructuring report for the full list of instances.
% - [ ] Fill in <REPO-URL-TODO> and <HASH-TODO> in
%   \S\ref{sec:evaluation-methodology} before submission.
% - [ ] x86 Linux throughput replication + RocksDB integration demo.
% - [ ] Author/affiliation + AI-assistance disclosure; fresh scoop-scan
%   <=2 weeks before submission.
% - [ ] Re-verify page/word count against the final DCC template if/when
%   the organizers publish one (current geometry is best-effort per
%   docs/DCC_STYLE.md's 6in x 9in text-area estimate).
